In this section we will describe some concept classes that can be learned under
the KWIK model.

\paragraph{Algorithm 1: Memorization Algorithm}
Given a finite input space $X$ and noise-free observations, any concept class
can be learned under the KWIK model with a KWIK bound of $|X|$ by using the
memorization algorithm. Which, for every received example $x \in X$:
\begin{enumerate}
  \item If it is the first time it receives $x$, the algorithm returns $\bot$,
  receives the correct output of $x$, and memorizes it.
  \item If it is not the first time it receives $x$, the algorithm has already
  memorized its correct label, so it outputs it.
\end{enumerate}

This algorithm will output $\bot$ at most once per example, because it is
impossible to see an example twice for the first time. Therefore, the KWIK bound
of the memorization algorithm is $|X|$. \\

\paragraph{Algorithm 2: Enumeration Algorithm}
Given any input space $X$ and noise-free observations, any finite concept class $C \subseteq X \to Y$
can be learned under the KWIK model with a KWIK bound of $|C| - 1$ by using the
elimination algorithm. The elimination algorithm maintains a set of ``potential
target concepts'' $\hat{H}$, initialized as $C$. For every received example
$x \in X$, it creates a set $\hat{L} = \{ h(x) | h \in \hat{H} \}$ and:
\begin{enumerate}
  \item If $|\hat{L}| > 1$, it outputs $\bot$ because at least 2 hypotheses
  disagree in predicting the output for $x$. When the correct output of $x$ is received, it
  updates $\hat{H}$ by eliminating from it all the hypotheses that
  incorrectly predicted the output for $x$.
  \item If $|\hat{L}| = 1$, it outputs the element in $\hat{L}$.
  \item If $|\hat{L}| = 0$, it means that the target concept is not in
  the concept class being learned.
\end{enumerate}

For every $\bot$ that the algorithm returns, it will eliminate at least one
concept from $\hat{H}$. Initially, $|\hat{H}| = |C|$. If the target
concept belongs to $C$, at the end of its execution, the algorithm will have
eliminated $|C| - 1$ concepts from $\hat{H}$. Therefore, the maximum number
of $\bot$s that the algorithm will output is $|C| - 1$.

\paragraph{Example 1:} The PhD lounge is frequented by a set $S$ of $n$ CLT
students. Every two weeks, any subset of them $G \subseteq S$ meets there to work
together on an assignment from the class. One of the students $f \in S$ is an
expert in slowing down the entire group, and whenever he is in the PhD lounge
($f \in G$), all the students in $G$ are forced to ``pull an all-nighter''.
Fortunately, there is another student $p \in S$ who is a genius, and, whenever
he is in the PhD lounge ($p \in G$), everyone is able to finish on time and go
home early, independently on the presence/absence of $f$. Rocco heard about the
story, and he wants to predict whether a subset of the students will ``pull an
all-nighter'' or not. However, he doesn't know the identity of $f$ and $p$. \\

The input space $X$ is all the possible subsets of $S$ ($|X| = 2^n$). Moreover,
the concept class to learn are all the possible assignments of $p$ and $f$
($|C| = n(n - 1)$). Therefore:
\begin{itemize}
  \item The memorization algorithm would achieve a KWIK bound of $2^n$.
  \item The enumeration algorithm would achieve a KWIK bound of $n(n - 1) - 1$.
\end{itemize}

\begin{figure}
  \centering
  \includegraphics[width=0.5\textwidth]{figures/algorithm3.pdf}
  \caption{Visualization of an execution of planar-distance Algorithm}
  \label{fig:sec4.1}
\end{figure}

\paragraph{Algorithm 3: Planar-distance Algorithm}
The two algorithms described above are limited to learning a finite concept
class $C$ or receiving a finite input space $X$. In this algorithm, this
restriction doesn't exist. \\

Let:
\begin{itemize}
  \item $ X = \mathbb{R}^2 $
  \item $ Y = \mathbb{R} $
  \item $ C = \{ f \ | \ f : X \to Y, \overline{x} \in \mathbb{R}^2, f(x) = \norm{x - \overline{x}} \} $
\end{itemize}

There is a point $\overline{x}$ unknown to the algorithm.  The target concept $c(x)$ fixes $\overline{x}$ and it outputs
the distance between an example $x$ and point $\overline{x}$. \\
\begin{enumerate}
  \item For an initial input $x_1$, the algorithm returns $\bot$ and it receives
  the distance $y_1$. Now, the algorithm knows that $\overline{x}$ is in the
  circumference $c_1$, with radius $y_1$ centered at $x_1$.
  \item When the algorithm receives the second example $x_2$:
  \begin{itemize}
    \item If $x_2 = x_1$, it returns $y_1$.
    \item Otherwise, it returns $\bot$ and receives the output $y_2$. Similarly,
    the algorithm knows that $\overline{x}$ lies in the circumference $c_2$ with
    radius $y_2$ centered at $x_2$. Then, $\overline{x}$ is restricted to one of
    the two intersecting points ($i_1$ and $i_2$) of $c_1$ and $c_2$, as illustrated in Figure
    \ref{fig:sec4.1} (b). At this point, if $c_1$ and $c_2$ are tangent, $i_1 =
    i_2 = \overline{x}$.
  \end{itemize}
  \item When the algorithm receives the third example $x_3$:
  \begin{itemize}
    \item If $x_3$ is collinear with $x_1$ and $x_2$, then, $\norm{x_3 - i_1} =
    \norm{x_3 - i_2}$. So, it returns $\norm{x_3 - i_1}$.
    \item If it is not collinear, then it returns $\bot$ and receives the output
    $y_3$. Similarly, the algorithm knows that $\overline{x}$ lies in the
    circumference $c_3$ with radius $y_3$ centered at $x_3$. Now, circumferences
    $c_1$, $c_2$ and $c_3$ only have one intersecting point, which is
    $\overline{x}$.
  \end{itemize}
\end{enumerate}

We just showed that, for any input set, the algorithm will completely identify
$\overline{x}$ by outputting at most $3$ $\bot$s. Therefore, the KWIK bound of the
planar-distance algorithm is $3$. As we saw in the case of $x_3$ being collinear
with $x_1$ and $x_2$, there are instances in which a KWIK algorithm can
have enough information for returning valid answers before having learned the
target concept.

\paragraph{Algorithm 4: Coin-Learning Algorithm}
We will now consider a simple KWIK algorithm which learns the probability of a
biased coin to come up heads. This algorithm will receive Bernoulli observations
($0$ when the coin comes up tails and $1$ when it comes up heads) and, with a
KWIK bound of $B(\varepsilon, \delta) = O(\frac{1}{\varepsilon^2}\log \frac{1}{\delta})$, it learns the
bias $p$ of the coin with $\varepsilon$-accuracy with probability at least $1-\delta$. \\

Let $\hat{p}$ be the estimate of $p$ made by the algorithm. It wants to estimate
the probability so that $|\hat{p} - p| \leq \varepsilon$. \\


Let $T$ be the number of observations it needs in order to satisfy the $\varepsilon$ and $\delta$ requirements.
Let $z_i$ be the $i$-th observation received by the algorithm.

\begin{itemize}
  \item If $i < T$ the algorithm can't guarantee that it will
  output an $\varepsilon$-accurate estimate with confidence $1-\delta$.
  \item Otherwise it calculates and outputs $\hat{p} = \frac{1}{T}\sum_{t=1}^{T}z_t$.
\end{itemize}
$T$ can be calculated by using a Hoeffding bound
\begin{eqnarray*}
  Pr[\hat{p} - p \geq \varepsilon] &\leq&  e^{-2T\varepsilon^2} \leq \delta/2 \\
  e^{-2T\varepsilon^2} &\leq& \delta/2 \\
  \Rightarrow T &\geq& \frac{1}{2\varepsilon^2}\ln \frac{2}{\delta}
\end{eqnarray*}

\begin{eqnarray*}
  Pr[\hat{p} - p \leq -\varepsilon] &\leq&  e^{-2T\varepsilon^2} \leq \delta/2 \\
  e^{-2T\varepsilon^2} &\leq& \delta/2 \\
  \Rightarrow T &\geq& \frac{1}{2\varepsilon^2}\ln \frac{2}{\delta}
\end{eqnarray*}

If $T = \frac{1}{2\varepsilon^2}\ln \frac{2}{\delta}$, with probability at least
$1 - \delta$, the algorithm will return an estimate $\hat{p}$ such that
$|\hat{p} - p| \leq \varepsilon$.

\paragraph{Algorithm 5: Noisy Linear-Regression Algorithm}
In section 3, we saw an algorithm for learning a weight vector $w$ such that
$f(x) = w \cdot x \ \forall x \in X$. In that case, what we observed was the
true value of $w \cdot x$. However, in the presence of additive noise, we don't
observe the true value, instead, we receive $w \cdot x + \eta$, with $\eta$
being the noise. This problem was solved by Stehl and Littman (2008)
\cite{Strehl}. The KWIK bound for this algorithm is $B(\varepsilon, \delta) =
O\left( \frac{d^3}{\varepsilon^4} \right)$, where $w \in \mathbb{R}^d$.
